% ==================================================================
% OUTLINE: Strong-Field and Full General Relativity
% from the Lattice State Packet in Conscious Point Physics
% Companion Paper to "Mechanistic Derivation of Relativistic Effects
% via Space Stress Vector (SSV) in the Dipole Sea" (Version 16)
% GR Companion II — Outline for Author Review
% 16 March 2026
% ==================================================================

\documentclass[12pt]{article}

\usepackage{amsmath,amssymb,amsthm}
\usepackage{geometry}
\usepackage{hyperref}
\usepackage{enumitem}
\usepackage{parskip}
\usepackage{xcolor}

\geometry{letterpaper, margin=1in}

\newcommand{\SSV}{\text{SSV}}
\newcommand{\PSR}{\text{PSR}}
\newcommand{\lP}{l_P}
\newcommand{\tP}{t_P}
\newcommand{\EP}{E_P}
\newcommand{\mP}{m_P}
\newcommand{\rS}{r_S}

% Annotation colours
\newcommand{\rigorous}[1]{\textcolor{blue}{#1}}
\newcommand{\qualitative}[1]{\textcolor{orange}{#1}}
\newcommand{\needsgrok}[1]{\textcolor{red}{#1}}
\newcommand{\deferred}[1]{\textcolor{gray}{#1}}

\title{\textbf{OUTLINE}\\[4pt]
Strong-Field and Full General Relativity\\
from the Lattice State Packet in Conscious Point Physics\\[6pt]
{\large GR Companion II ---
Companion Paper to ``Mechanistic Derivation of Relativistic Effects\\
via Space Stress Vector (SSV) in the Dipole Sea'' (Version~16)}}

\author{Thomas Lee Abshier, ND\\
Hyperphysics Institute\\
\texttt{https://hyperphysics.com}}

\date{16 March 2026 --- Pre-writing outline}

\begin{document}
\maketitle

\noindent\textbf{Legend:}
\rigorous{Blue = rigorous, can write now.}\quad
\qualitative{Orange = qualitative / sketch, needs explicit derivation.}\quad
\needsgrok{Red = needs Grok input before writing.}\quad
\deferred{Gray = deferred to a later companion.}

%======================================================================
\section*{Paper title and scope}
%======================================================================

\textbf{Proposed title:} ``Strong-Field General Relativity and the CPP
Field Equation from the Lattice State Packet''

\textbf{Scope:} This paper extends GR Companion~I (weak-field) to:
\begin{itemize}[noitemsep]
  \item The exact (nonlinear) Schwarzschild solution from the full PSR
    formula, without linearisation.
  \item The CPP self-consistency condition as the discrete-lattice
    analog of the Einstein field equations.
  \item Resolution of the Schwarzschild singularity via the CP
    Exclusion Rule.
  \item Gravitational redshift and light deflection at all orders
    (not just first order).
  \item Kerr metric from rotational $\mathbf{\SSV}_{\rm net}$ (qualitative;
    Grok input needed).
  \item Cosmological implications: the CPP vacuum SSV as a candidate
    for the cosmological constant.
\end{itemize}

The paper does \textbf{not} attempt:
\begin{itemize}[noitemsep]
  \item Full discrete-to-continuum convergence proof (this is a
    dedicated future companion — requires 600-cell spectral analysis
    analogous to Spin~III).
  \item Quantum gravity / Planck-scale foam (out of scope for this series).
\end{itemize}

%======================================================================
\section*{Plain Language Summary (draft)}
%======================================================================

GR Companion~I derived the Schwarzschild metric in the weak-field limit:
close to a mass but far from its surface.  This companion goes further
--- to the edge of a black hole, to the Planck-scale resolution of the
singularity, and to the general rule that replaces Einstein's field
equations.

Near a black hole, the gravitational SSV becomes so intense that the
PSR --- the spatial budget each CP has for its next hop --- contracts
nearly to zero.  At the Schwarzschild radius the PSR reaches exactly
half the Planck length, and the GP Exclusion Rule prevents any further
contraction.  There is no singularity in CPP: the black hole interior
is a region of maximally compressed lattice, not a point of infinite
density.

The rule that governs how mass-energy shapes the lattice --- the CPP
field equation --- is a self-consistency condition: the metric you
compute from the LSP broadcasts must be the same metric that governs
how LSP broadcasts propagate.  In the weak-field limit this reduces
exactly to Einstein's linearised field equations.  This paper shows
the strong-field version.

%======================================================================
\section*{Section-by-section outline}
%======================================================================

\medskip
\noindent\rule{\textwidth}{0.4pt}

\subsection*{Section 1 — Introduction}

\begin{itemize}[noitemsep]
  \item \rigorous{Summarise GR~I results: weak-field Schwarzschild,
    Eddington factor-of-two, geodesic equation, gravitational waves,
    effective horizon.}
  \item \rigorous{State the three open problems from GR~I that this paper
    closes: (1) exact Schwarzschild, (2) full field equation, (3)
    singularity resolution.}
  \item \rigorous{State Kerr as a fourth result — rigorous at the
    mechanism level, explicit calculation deferred.}
  \item \rigorous{Preview: the PSR formula is already exact; the
    weak-field paper only used its linearised limit.  The full
    nonlinear PSR formula gives exact GR without new postulates.}
\end{itemize}

\medskip\noindent\rule{\textwidth}{0.4pt}

\subsection*{Section 2 — The Full Nonlinear PSR Formula}

\textit{Key result: the PSR formula is already exact GR --- the
weak-field paper only used it to first order.}

\begin{itemize}[noitemsep]
  \item \rigorous{Restate the PSR formula from the main paper:
    \[ \PSR_{\rm eff} = \frac{\lP}{1 + k\cdot\Delta|\SSV|_{\rm abs}} \]}
  \item \rigorous{The metric component $g_{tt}$ from GR~I is:
    \[ g_{tt} = \PSR_{\rm eff}/\lP = \frac{1}{1+k\Delta|\SSV|} \]}
  \item \needsgrok{At what value of $k\Delta|\SSV|$ does this equal the
    exact Schwarzschild $g_{tt} = 1 - r_S/r$?  Need to confirm whether
    the CPP formula gives the Schwarzschild result exactly or only
    to first order.}
  \item \rigorous{At the Schwarzschild radius $r_S = 2GM/c^2$:
    $k\Delta|\SSV| = 1$ (from GR~I, Eq.~\ref{eq:horizon}), giving
    $\PSR_{\rm eff} = \lP/2$.}
  \item \rigorous{For all $r > r_S$: $\PSR_{\rm eff}$ is positive,
    finite, and analytic. No singularity.}
  \item \rigorous{For $r \leq r_S$: CP Exclusion prevents $\PSR < \lP/2$,
    giving a natural ultraviolet cutoff.}
\end{itemize}

\textbf{What Grok should check:} Does the CPP metric
$g_{tt} = 1/(1 + GM/rc^2 \cdot k^{-1})$ reduce to the exact
Schwarzschild $g_{tt} = 1 - r_S/r$ by a coordinate identification,
or only approximately?  The SSV source term from GR~I has
$k\Delta|\SSV| = GM/(rc^2)$, giving $g_{tt} = 1/(1 + GM/rc^2)$,
which equals $1 - GM/rc^2$ to first order but is \emph{not} exact
Schwarzschild.  The resolution may require the isotropic Schwarzschild
coordinate form rather than the standard form.

\medskip\noindent\rule{\textwidth}{0.4pt}

\subsection*{Section 3 — Exact Schwarzschild from the Self-Consistency
Condition}

\textit{Key result: the exact metric follows from requiring LSP
propagation to be consistent with the metric it generates.}

\begin{itemize}[noitemsep]
  \item \rigorous{Define the CPP self-consistency condition precisely:
    the metric $g_{\mu\nu}$ reconstructed from LSP summation must equal
    the metric that governs LSP propagation speed and direction.}
  \item \qualitative{In the static, spherically symmetric case, show
    that the self-consistency condition is equivalent to the
    Schwarzschild ODE:
    \[ \frac{d}{dr}\left[r^2\frac{d\Phi}{dr}\right] = 0 \]
    outside the mass, with $\Phi$ the gravitational potential sourcing
    $\Delta|\SSV|$.}
  \item \rigorous{Exact Schwarzschild metric in isotropic coordinates:
    \[ ds^2 = -\left(\frac{1-\rho}{1+\rho}\right)^2 c^2\,dt^2
       + \left(1+\rho\right)^4 (d\mathbf{x})^2, \quad
       \rho = \frac{GM}{2c^2 r} \]}
  \item \needsgrok{Show that the CPP metric (both $g_{tt}$ from
    $|\SSV|_{\rm abs}$ and $g_{ij}$ from $\SSV_{\rm net}$) reproduces
    this isotropic form exactly, not just to first order.}
  \item \rigorous{Derive the exact gravitational redshift:
    \[ \frac{\nu_{\rm obs}}{\nu_{\rm em}}
       = \sqrt{\frac{g_{tt}(r_{\rm em})}{g_{tt}(r_{\rm obs})}}
       = \frac{1 - r_S/2r_{\rm obs}}{1 - r_S/2r_{\rm em}} \]
    from the LSP frequency broadcast formula.}
\end{itemize}

\medskip\noindent\rule{\textwidth}{0.4pt}

\subsection*{Section 4 — The CPP Field Equation}

\textit{Key result: the self-consistency condition, when written as a
differential equation, is the discrete-lattice analog of the Einstein
field equations.}

\begin{itemize}[noitemsep]
  \item \rigorous{State the self-consistency condition as a fixed-point
    equation: $\mathcal{M}[\rho_{\rm LSP}] = \mathcal{G}[\mathcal{M}]$,
    where $\mathcal{M}$ is the metric reconstruction map and
    $\mathcal{G}$ is the geodesic propagation operator.}
  \item \qualitative{In the continuum limit, expand the discrete LSP
    sum in powers of $\lP/r$ and show the leading term gives:
    \[ G_{\mu\nu} = \frac{8\pi G}{c^4} T_{\mu\nu} \]
    where $T_{\mu\nu}$ is the LSP energy-momentum content.}
  \item \rigorous{The weak-field limit is already established (GR~I,
    Section~5): $\Box\bar{h}_{\mu\nu} = -16\pi G/c^4 \cdot T_{\mu\nu}$.
    This section establishes the strong-field generalization.}
  \item \deferred{Full convergence proof (discrete sum to smooth
    Riemannian manifold): deferred to GR Companion~III, using the
    same 24-cell Voronoi machinery developed in Spin~III.}
  \item \rigorous{State the CPP prediction: the Einstein field
    equations hold \emph{exactly} in the continuum limit, with
    corrections of order $(\lP/r)^2$ --- negligible at all scales
    above the Planck length.}
\end{itemize}

\medskip\noindent\rule{\textwidth}{0.4pt}

\subsection*{Section 5 — Singularity Resolution via CP Exclusion}

\textit{Key result: the Schwarzschild singularity is replaced by a
Planck-density core.  No new physics required.}

\begin{itemize}[noitemsep]
  \item \rigorous{At $r = r_S$: $\PSR_{\rm eff} = \lP/2$ (from GR~I).
    Established. Restate without reproof.}
  \item \rigorous{For $r < r_S$: the self-consistency condition would
    require $\PSR_{\rm eff} < \lP/2$.  The CP Exclusion Rule
    (one CP per GP per tick) forbids this.}
  \item \rigorous{Consequence: infalling CPs are deflected to
    neighboring GPs rather than annihilated.  The black hole interior
    is a region of \emph{maximally compressed lattice} with
    $\PSR = \lP/2$ everywhere, not a point singularity.}
  \item \qualitative{The interior metric:
    \[ ds^2_{\rm interior} \approx -\left(\frac{\lP/2}{r}\right)^2 c^2\,dt^2
       + \left(\frac{r}{\lP/2}\right)^2 (d\mathbf{x})^2 \]
    (de~Sitter-like near the core).  This is a conjecture; explicit
    derivation needed.}
  \item \rigorous{The information paradox: LSP information is
    preserved in the compressed-lattice interior.  CPs are deflected,
    not destroyed.  The unitarity of CP state evolution is maintained
    by construction.}
  \item \rigorous{Remark on Hawking radiation: thermal LSP emission
    from the $\PSR = \lP/2$ boundary is the CPP mechanism.  The
    Hawking temperature $T_H = \hbar c^3/(8\pi G M k_B)$ is expected
    to follow from the ZBW frequency at the PSR boundary.
    Explicit derivation deferred.}
\end{itemize}

\medskip\noindent\rule{\textwidth}{0.4pt}

\subsection*{Section 6 — Exact Light Deflection and Perihelion Precession}

\textit{Key result: all-orders deflection from the exact CPP metric.}

\begin{itemize}[noitemsep]
  \item \rigorous{GR~I established the first-order (Eddington)
    deflection $\delta\phi = 4GM/bc^2 = 1.750''$ from the
    two-component LSP.  This section generalises to all orders.}
  \item \rigorous{From the exact Schwarzschild (or isotropic)
    metric, the deflection integral:
    \[ \Delta\phi = \int_{-\infty}^{\infty}
       \frac{d}{dr}\left(\frac{n_{\rm eff}(r)}{c}\right)
       \frac{b\,dr}{\sqrt{r^2 - b^2}} \]
    where $n_{\rm eff}(r) = c/(\PSR_{\rm eff}/\lP)$ is the
    effective refractive index of the CPP lattice near mass $M$.}
  \item \qualitative{Higher-order terms reproduce the GR post-Newtonian
    series; CPP predicts the same coefficients as GR.}
  \item \rigorous{Mercury perihelion precession: from the exact
    geodesic equation (established in GR~I, Section~4) applied to
    the full Schwarzschild metric, the precession per orbit is
    $\Delta\phi_{\rm prec} = 6\pi GM/(c^2 a(1-e^2))$.
    The CPP derivation follows the standard GR calculation with
    the PSR formula as the substitute for the metric.}
\end{itemize}

\medskip\noindent\rule{\textwidth}{0.4pt}

\subsection*{Section 7 — Kerr Metric from Rotational SSV}

\textit{Key result (mechanism): frame dragging from rotational
$\mathbf{\SSV}_{\rm net}$ curl.}

\begin{itemize}[noitemsep]
  \item \rigorous{The Kerr metric requires an off-diagonal component
    $g_{t\phi} \neq 0$.  In CPP, this arises from a curl in
    $\mathbf{\SSV}_{\rm net}$ generated by the rotating mass.}
  \item \qualitative{A rotating mass imparts angular momentum to the
    local DP sea.  The DP sea rotation produces a circulating
    $\mathbf{\SSV}_{\rm net}$ with curl proportional to the angular
    momentum $J$ of the source.}
  \item \needsgrok{Explicit calculation: does the rotational
    $\mathbf{\SSV}_{\rm net}$ curl produce the correct Kerr
    $g_{t\phi} = -2GJr\sin^2\theta/c^2 r^3$ in the weak-field limit?
    This is the central question for Grok.}
  \item \qualitative{Frame dragging (Lense-Thirring effect) in CPP:
    a gyroscope precesses because the rotating $\mathbf{\SSV}_{\rm net}$
    preferentially selects oriented GP edges for the gyroscope's ZBW
    oscillation.}
  \item \deferred{Kerr-Newman metric (charged rotating black hole):
    combines rotational SSV curl with charge-polarity SSV.
    Mechanism clear; derivation deferred.}
\end{itemize}

\medskip\noindent\rule{\textwidth}{0.4pt}

\subsection*{Section 8 — The Cosmological Constant from Vacuum SSV}

\textit{Key result: the residual SSV of the unperturbed DP sea is a
candidate for $\Lambda$.}

\begin{itemize}[noitemsep]
  \item \rigorous{The vacuum DP sea has a non-zero background
    $|\SSV|_{\rm abs} = |\SSV|_{\rm vac}$ set by the 600-cell geometry
    (established in the Stiffness~$C$ companion).}
  \item \qualitative{This uniform background SSV acts as a
    cosmological constant:
    \[ \Lambda_{\rm CPP} = k\cdot|\SSV|_{\rm vac} \cdot \frac{3}{c^2}. \]}
  \item \needsgrok{What is the numerical value of $|\SSV|_{\rm vac}$
    from the 600-cell geometry?  Does it reproduce the observed
    $\Lambda \approx 1.1\times10^{-52}\ \text{m}^{-2}$?}
  \item \rigorous{The CPP cosmological constant does not suffer from
    the $10^{120}$ fine-tuning problem of QFT: the vacuum energy is
    not the zero-point energy of all quantum fields but the
    specific SSV of the DP sea, which is bounded by the 600-cell
    geometry.}
\end{itemize}

\medskip\noindent\rule{\textwidth}{0.4pt}

\subsection*{Section 9 — Consistency with GR Companion I and the SR Series}

\begin{itemize}[noitemsep]
  \item \rigorous{Every result of GR~I is recovered as the
    $k\Delta|\SSV| \ll 1$ limit.}
  \item \rigorous{The Eddington factor-of-two lensing prediction is
    unchanged.}
  \item \rigorous{The SR companion results are the flat-space
    ($\Delta|\SSV| = 0$) limit.}
  \item \rigorous{No new postulates introduced beyond those of GR~I.}
\end{itemize}

\medskip\noindent\rule{\textwidth}{0.4pt}

\subsection*{Section 10 — Open Problems}

\begin{enumerate}[noitemsep]
  \item \textbf{Discrete-to-continuum convergence} (GR Companion~III):
    proving that the discrete 600-cell lattice LSP sum converges to the
    smooth Riemannian manifold, with corrections of order $(\lP/r)^2$.
    Uses the 24-cell Voronoi eigenvalue machinery from Spin~III.
  \item \textbf{Kerr explicit calculation}: the rotational
    $\mathbf{\SSV}_{\rm net}$ derivation of $g_{t\phi}$ (needs Grok).
  \item \textbf{Hawking radiation}: explicit derivation of $T_H$ from
    ZBW emission at the $\PSR = \lP/2$ boundary.
  \item \textbf{Cosmological constant numerical match}: computing
    $|\SSV|_{\rm vac}$ from 600-cell geometry and comparing to
    observation.
  \item \textbf{Big Bang / GP Exclusion cosmology}: initial Moment
    CP crowding drives expansion without an inflaton.
    Treated in a dedicated cosmology companion.
\end{enumerate}

\medskip\noindent\rule{\textwidth}{0.4pt}

\subsection*{Section 11 — Conclusion}

Rigorous results (to be listed in numbered form):
\begin{enumerate}[noitemsep]
  \item Full nonlinear PSR formula as the exact metric.
  \item CPP self-consistency condition as the lattice field equation.
  \item Singularity resolution via CP Exclusion (Planck-density core).
  \item Exact gravitational redshift from LSP frequency broadcast.
  \item All-orders light deflection and perihelion precession.
  \item Kerr mechanism (qualitative; calculation deferred).
  \item CPP vacuum SSV as cosmological constant candidate (qualitative).
\end{enumerate}

%======================================================================
\section*{Questions for Grok before writing begins}
%======================================================================

\begin{enumerate}
  \item \textbf{[Critical]} Does $g_{tt} = 1/(1 + GM/rc^2)$ from the
    CPP PSR formula equal the exact Schwarzschild
    $g_{tt} = 1 - r_S/r$ by a coordinate identification?
    Specifically: does the isotropic Schwarzschild coordinate form
    close the gap?

  \item \textbf{[Critical]} What is the CPP self-consistency condition
    written as an explicit equation?  The Newtonian Gravity companion
    describes it as a feedback loop qualitatively.  Has it been
    written as $\mathcal{F}[g_{\mu\nu}, T_{\mu\nu}] = 0$?

  \item \textbf{[Important]} Does the rotational $\mathbf{\SSV}_{\rm net}$
    curl from a spinning mass reproduce the Kerr
    $g_{t\phi} = -2GJr\sin^2\theta/(c^2 r^3)$ in the weak-field limit?

  \item \textbf{[Useful]} What is the current CPP estimate for the
    vacuum SSV $|\SSV|_{\rm vac}$ from the 600-cell geometry?

  \item \textbf{[Useful]} Is there existing Grok material on the
    CPP black hole interior / de~Sitter-like core?
\end{enumerate}

%======================================================================
\section*{Paper length and timeline estimate}
%======================================================================

\begin{itemize}[noitemsep]
  \item Sections 1--3 (intro, nonlinear PSR, exact Schwarzschild):
    writable now if Grok answers Question~1 affirmatively.
  \item Section~4 (CPP field equation): writable if Grok answers
    Question~2.  Otherwise stated as strong conjecture with
    the weak-field limit as the rigorous anchor.
  \item Section~5 (singularity resolution): writable now ---
    the CP Exclusion argument is already in GR~I.
  \item Section~6 (exact deflection, perihelion): writable now
    from the standard GR calculation with PSR substitution.
  \item Section~7 (Kerr): writable at mechanism level now;
    explicit calculation after Grok Question~3.
  \item Section~8 (cosmological constant): writable at mechanism
    level now; numerical match after Grok Question~4.
  \item Sections 9--11: writable now.
  \item \textbf{Estimated length:} 18--22 pages (comparable to GR~I).
  \item \textbf{Estimated writing time:} one session after Grok
    answers Questions 1 and 2.
\end{itemize}

\end{document}
